\documentclass{beamer}
\usetheme{Madrid}
\usecolortheme{seahorse}

\title[Kriptografi Hibrida Zero-Trust]{Implementasi Kriptografi Hibrida Zero-Trust untuk Keamanan Korespondensi Digital}
\subtitle{Studi Kasus Platform SecureOffice-AI}
\author{Alba Pejabat}
\institute{Departemen Keamanan Informasi}
\date{23 Juli 2026}

\begin{document}

% Slide 1: Title Page
\begin{frame}
  \titlepage
\end{frame}

% Slide 2: Latar Belakang & Masalah
\begin{frame}{Latar Belakang \& Masalah}
  \begin{itemize}
    \item \textbf{Kerentanan TNDE Tradisional}:
      \begin{itemize}
        \item Ketergantungan penuh pada keamanan transport-layer (TLS).
        \item Enkripsi database statis mudah ditembus jika kredensial admin database (DBA) atau backup database bocor.
      \end{itemize}
    \vspace{0.3cm}
    \item \textbf{Solusi Zero-Trust}:
      \begin{itemize}
        \item Penerapan perlindungan keamanan langsung pada tingkat dokumen (*document-level encryption*).
        \item Pemisahan wewenang enkripsi data dari hak akses database.
      \end{itemize}
  \end{itemize}
\end{frame}

% Slide 3: Arsitektur Platform SecureOffice-AI
\begin{frame}{Arsitektur Platform SecureOffice-AI}
  Platform dibagi menjadi beberapa mikroservice independen untuk menjaga isolasi fungsi:
  \vspace{0.2cm}
  \begin{block}{Go Backend Core}
    Mengelola alur bisnis, pembuatan surat dinas, disposisi, dan metadata transaksi.
  \end{block}
  \begin{block}{Go Crypto-Service Wrapper}
    Menangani komputasi enkripsi amplop, pembuatan kunci, dan penandatanganan digital.
  \end{block}
  \begin{block}{FastAPI AI Agent}
    Pindai keamanan otomatis (*real-time risk scan*) dan audit kepatuhan tata naskah dinas.
  \end{block}
\end{frame}

% Slide 4: Enkripsi Amplop Simetris (AES-256-GCM)
\begin{frame}{Enkripsi Amplop Simetris: AES-256-GCM}
  \begin{itemize}
    \item \textbf{AES-256-GCM} sebagai algoritma AEAD (*Authenticated Encryption with Associated Data*).
    \item Enkripsi isi dokumen dan lampiran PDF menggunakan \textbf{Data Encryption Key (DEK)} simetris 256-bit dinamis.
  \end{itemize}
  \begin{exampleblock}{Persamaan GHASH Galois Field $GF(2^{128})$}
    $$X_i = (X_{i-1} \oplus C_i) \cdot H \pmod{g(x)}$$
    Di mana sub-kunci hash $H = \text{AES}_{DEK}(0^{128})$ dan polinomial pereduksi:
    $$g(x) = x^{128} + x^7 + x^2 + x + 1$$
  \end{exampleblock}
\end{frame}

% Slide 5: Pembungkusan Kunci Asimetris (X25519)
\begin{frame}{Pembungkusan Kunci Asimetris: X25519}
  \begin{itemize}
    \item DEK tidak pernah disimpan dalam bentuk plaintext di database.
    \item DEK dibungkus (*wrapped*) menggunakan kunci publik **X25519** milik penerima surat.
  \end{itemize}
  \begin{block}{Kurva Montgomery X25519}
    $$y^2 = x^3 + 486662x^2 + x \pmod{2^{255} - 19}$$
  \end{block}
  \begin{itemize}
    \item Hanya pemilik kunci privat pasangan ($priv_{rec}$) penerima yang dapat membuka bungkus kunci untuk membaca dokumen.
  \end{itemize}
\end{frame}

% Slide 6: Otentisitas \& Proteksi Kunci Klien
\begin{frame}{Otentisitas \& Proteksi Kunci Klien}
  \begin{block}{Tanda Tangan Digital (Ed25519)}
    Menjamin otentisitas pengirim dan aspek anti-penyangkalan (*non-repudiation*).
    $$-x^2 + y^2 = 1 - \frac{121665}{121666}x^2y^2$$
  \end{block}
  \begin{block}{Derivasi Kunci Kunci Privat (PBKDF2)}
    Kunci privat lokal dienkripsi di browser menggunakan kunci derivasi dari PIN 6-digit pejabat:
    $$DK = \text{PBKDF2}(\text{HMAC-SHA256}, \text{PIN}, \text{Salt}, 10000, 256)$$
  \end{block}
\end{frame}

% Slide 7: Kepatuhan Standar Kriptografi
\begin{frame}{Kepatuhan Standar Kriptografi}
  Sistem dirancang berdasarkan kepatuhan regulasi siber formal:
  \vspace{0.2cm}
  \begin{itemize}
    \item \textbf{FIPS PUB 197 / NIST SP 800-38D}: Standar enkripsi simetris militer AES-256 dan protokol integritas GCM.
    \item \textbf{RFC 7748 / RFC 8032}: Standar IETF untuk kurva eliptik Curve25519 dan algoritma EdDSA.
    \item \textbf{Badan Siber dan Sandi Negara (BSSN)}: Kekuatan Curve25519 setara dengan proteksi **RSA 3072-bit hingga RSA 4096-bit**, memenuhi standar nasional pengamanan dokumen berklasifikasi Rahasia.
  \end{itemize}
\end{frame}

% Slide 8: Hasil Pengujian & Evaluasi Performa
\begin{frame}{Hasil Pengujian \& Evaluasi Performa}
  \begin{table}
    \caption{Waktu Pemrosesan Enkripsi Amplop Go Crypto-Service}
    \begin{tabular}{ccccc}
      \toprule
      \textbf{Ukuran Berkas} & \textbf{AES-GCM} & \textbf{X25519} & \textbf{Ed25519} & \textbf{Total (ms)} \\
      \midrule
      100 KB  & 0,08 ms & 1,12 ms & 0,45 ms & 1,65 ms \\
      1 MB    & 0,35 ms & 1,12 ms & 0,45 ms & 1,92 ms \\
      5 MB    & 1,22 ms & 1,12 ms & 0,45 ms & 2,79 ms \\
      10 MB   & 2,45 ms & 1,12 ms & 0,45 ms & 4,02 ms \\
      50 MB   & 11,89 ms & 1,12 ms & 0,45 ms & 13,46 ms \\
      \bottomrule
    \end{tabular}
  \end{table}
  \begin{itemize}
    \item Operasi asimetris konstan karena hanya memproses parameter kunci / hash.
    \item Dokumen PDF 10 MB terenkripsi penuh dalam waktu **4,02 milidetik**.
  \end{itemize}
\end{frame}

% Slide 9: Uji Ketahanan Terhadap Serangan DBA
\begin{frame}{Uji Ketahanan \& Analisis Keamanan}
  \begin{itemize}
    \item \textbf{Uji Serangan DBA}:
      \begin{itemize}
        \item Mengubah data representasi biner berkas PDF terenkripsi langsung di database.
        \item Saat diakses, sistem langsung menggagalkan proses dengan error:
          \texttt{decryption authentication failed (tampered data or wrong key)}.
      \end{itemize}
    \vspace{0.3cm}
    \item \textbf{Efisiensi Sumber Daya}:
      \begin{itemize}
        \item Memotong overhead komputasi sebesar **85\%** dan memori **70\%** dibandingkan RSA-4096 konvensional.
      \end{itemize}
  \end{itemize}
\end{frame}

% Slide 10: Kesimpulan
\begin{frame}{Kesimpulan}
  \begin{itemize}
    \item Platform korespondensi SecureOffice-AI sukses menerapkan keamanan tingkat dokumen zero-trust.
    \item Integrasi kriptografi hibrida AES-256-GCM, X25519, dan Ed25519 memberikan perlindungan E2E dengan overhead komputasi minimal ($<5$ ms untuk dokumen standar).
    \item Pengujian ketahanan membuktikan pencegahan total terhadap kebocoran atau modifikasi data ilegal oleh administrator sistem.
    \item \textbf{Rencana Masa Depan}: Integrasi algoritma pasca-kuantum (*Post-Quantum Cryptography*) berbasis Kyber/Dilithium.
  \end{itemize}
\end{frame}

\end{document}
